\scanpage{108}
% archive sheet number: 107
\uline{Remarque.} Nous nous sommes bornés ici aux formules générales valables sans conditions de régularité (genre pseudo-cohérence, ou quasi-cohérence) qui donnent de nombreuses autres formules. D'autre part, nous nous sommes bornés aux cas d'existence \inserted{“naturelle” i.e. générale} des op. envisagées $\overset{\mathbb L}{\otimes}$, \inserted{$\mathbb R\mathscr H\!om$}, $\mathbb L f^*$, $\mathbb R f_*$, sans examiner les variantes \unsure{pr.|pro}\deleted{\unsure{venant|venant de}} d'hyp. de Tor-dim ou Ext-dim. finie \ldots.

\begin{center}\rule{0.95cm}{0.4pt}\\[-0.75em]\rule{0.72cm}{0.4pt}\end{center}

\textcircled{2} Formules \unsure{sur les précédentes|sur les précédents|pour les précédentes|pour les précédents}, faisant intervenir égalemt $!$ en plus de $\overset{\mathbb L}{\otimes}$ : Künneth.

\deleted{\uline{\unsure{Théorème}}. Soit $f:X\longrightarrow Y$ un morphisme compactifiable.}
\deleted{Nous \unsure{pouvons|pourr}}

Notons que en général nous sommes obligés de travailler avec des faisceaux d'anneaux constant $\mathbb Z/n\mathbb Z=A$ qui n'est pas un corps ni de tor-dimension \inserted{globale} finie, donc dès qu'on commence à faire du $\overset{\mathbb L}{\otimes}$, \deleted{\unsure{en}} en \unsure{partant|faisant} avec un faisceau réduit au degré 0, on va \unsure{tomber|déborder} \unsure{dans|des} $D^-$. Il faut
